%\section{Mindmap}
%\vfill
\begin{figure}[!t]
	\centering
	\begin{tikzpicture}[
			mindmap,
			text width=2.5cm,
			align=center,
			every node/.style={concept},
			concept color=gray!30,
			level 1/.append style={
				level distance=5cm,
				minimum size=3cm,
				text width=2.8cm,
				font=\large,
			},
			level 2/.append style={
				level distance=3.5cm,
				minimum size=2cm,
				text width=1.8cm
			},
			]
			
			\node[minimum size=4.5cm, text width=4.3cm, line width=6pt, fill=gray!10] (TTT) { \Large\textbf{Estimação com tempo de trânsito} }
			child[concept color=olive!10,grow=225] { 
				node {\textbf{Método}} 
				child[grow=120] { node {Estado adjunto \\(2)} }
				child[grow=180] { node {TTT \\(8)} }
				child[grow=240] { node {Redes neurais \\(6)} }
				child[grow=300] { node {Outros \\(6)} }
			}
			child[concept color=purple!10, grow=300] {
				node {\textbf{Aplicação}}
				child[grow=240] { node {Geologia \\(11)} }
				child[grow=300] { node {Termometria \\(5)} }
				child[grow=60] {node {Outros \\(3)} }
			}
			child[concept color=brown!10, grow=60] {
				node {\textbf{Dados de entrada}} 
				child[grow=150] { node {Tempo de trânsito \\(16)} }
				child[grow=0] { node {Modelo inicial \\(3)} }
				child[grow=-45] { node {Outros \\(3)}}
			}
			;	
	\end{tikzpicture}
	\caption[\textsl{Mindmap} da \cref{sec_rt}]{\textsl{Mindmap} da \cref{sec_rt} sobre inversão com tempo de trânsito.}
	\label{fig:mindmap_sec3}
\end{figure}